\section{Giới thiệu}
\subsection{Đặt vấn đề}

Trong bối cảnh giáo dục đại học tại Việt Nam, Trường Đại học Bách Khoa -- Đại học Quốc gia TP.HCM (ĐH Bách Khoa) là một trong những cơ sở đào tạo hàng đầu, với quy mô sinh viên lên đến hơn 27.000 người. Đây không chỉ là con số lớn nhất tại TP.HCM mà còn phản ánh vị thế dẫn đầu về số lượng đoàn viên và hội viên tham gia các hoạt động của Đoàn Thanh niên và Hội Sinh viên.  

Tuy nhiên, quy mô khổng lồ này lại đối mặt với một nghịch lý lớn: công tác quản lý hiện tại chủ yếu dựa vào các phương pháp thủ công và rời rạc, nguyên nhân là do đoàn hội hiện nay vẫn chưa có hệ thống quản lý thông tin tập trung.

Cụ thể, toàn bộ dữ liệu liên quan đến người dùng (đoàn viên, hội viên), hoạt động và tài chính không được lưu trữ trên một máy chủ thống nhất mà phân tán trong tài liệu giấy được lưu trong kho và các tệp Word, Excel trên máy tính cá nhân của cán bộ. Việc tra cứu, tổng hợp báo cáo hay liên kết dữ liệu đòi hỏi nhiều thao tác thủ công, gây tốn kém thời gian và nhân lực. Sinh viên phải ``bơi'' trong biển thông tin trên mạng xã hội để tìm kiếm hoạt động phù hợp, trong khi cán bộ gặp nhiều khó khăn trong công tác quản lý và thống kê.

Từ thực tế trên, thông qua quá trình khảo sát tại Đoàn trường và Hội Sinh viên, có thể xác định ba vấn đề cốt lõi cần được giải quyết:

\begin{itemize}
    \item \textbf{Quản lý Người dùng và Tổ chức:}  
    Chưa tồn tại hệ thống định danh số hóa cho đoàn viên và hội viên. Việc chứng minh tư cách chủ yếu dựa trên sổ đoàn hoặc hồ sơ giấy tờ thủ công, dễ thất lạc hoặc hư hỏng, dẫn đến việc xác minh thông tin tốn nhiều thời gian và nhân lực (ví dụ: kiểm tra sổ đoàn gốc hoặc biên lai đóng hội phí).

    \item \textbf{Quản lý Hoạt động:}  
    Thông tin về các chương trình, phong trào chủ yếu được đăng tải rải rác trên fanpage hoặc mạng xã hội, gây ra tình trạng trùng lặp, chồng chéo dữ liệu và thiếu công cụ lọc, thống kê tự động. Sinh viên phải tự tìm kiếm thông tin, trong khi cán bộ khó theo dõi và tổng hợp số lượng hoạt động theo từng giai đoạn.

    \item \textbf{Quản lý Tài chính:}  
    Quy trình thu -- chi chủ yếu dựa trên tiền mặt, với dòng tiền đi từ dưới lên (lớp thu, nộp khoa; khoa tổng hợp, nộp trường). Cách thức này tiềm ẩn nhiều rủi ro như thất thoát, sai lệch số liệu, thiếu minh bạch, đồng thời tạo áp lực và trách nhiệm lớn cho các cán bộ trung gian.
\end{itemize}

Những hạn chế nêu trên không chỉ làm giảm hiệu quả hoạt động của Đoàn -- Hội mà còn ảnh hưởng trực tiếp đến quyền lợi và mức độ tham gia của sinh viên. Do đó, việc xây dựng một giải pháp chuyển đổi số toàn diện là yêu cầu cấp thiết nhằm nâng cao tính minh bạch, hiệu quả quản lý và khả năng mở rộng trong tương lai.
\subsection{Các hướng giải quyết liên quan}

Xuất phát từ các vấn đề cốt lõi đã được phân tích, có thể nhận thấy rằng những hạn chế trong công tác quản lý Đoàn -- Hội không thể được khắc phục bằng các giải pháp rời rạc hay mang tính cục bộ. Thay vào đó, cần một hướng tiếp cận tổng thể dựa trên nền tảng chuyển đổi số, nhằm xây dựng một hệ thống quản lý thống nhất, đồng bộ và có khả năng mở rộng trong tương lai.

Đối với vấn đề quản lý người dùng và tổ chức, hướng giải quyết trọng tâm là số hóa dữ liệu định danh của Đoàn viên và Hội viên. Việc xây dựng hồ sơ điện tử cho mỗi sinh viên cho phép lưu trữ tập trung các thông tin về tư cách Đoàn viên, Hội viên cũng như lịch sử tham gia hoạt động. Qua đó, quá trình xác minh thông tin được thực hiện nhanh chóng, chính xác và giảm thiểu sự phụ thuộc vào hồ sơ giấy truyền thống.

Đối với vấn đề quản lý hoạt động, hướng giải quyết được đặt ra là tập trung hóa và chuẩn hóa thông tin hoạt động trên một nền tảng duy nhất. Các chương trình, phong trào và hoạt động xã hội sẽ được đăng tải, quản lý và theo dõi tập trung, giúp sinh viên dễ dàng tiếp cận và đăng ký tham gia, đồng thời hỗ trợ cán bộ trong việc thống kê, theo dõi và đánh giá hiệu quả tổ chức hoạt động theo từng giai đoạn.

Đối với vấn đề quản lý tài chính, hướng giải quyết chủ yếu là số hóa quy trình thu -- chi và từng bước giảm sự phụ thuộc vào tiền mặt. Việc quản lý các khoản thu, chi thông qua hệ thống điện tử giúp đảm bảo tính minh bạch, khả năng truy xuất và kiểm soát dòng tiền, đồng thời hạn chế các rủi ro liên quan đến thất thoát và sai lệch số liệu. Bên cạnh đó, giải pháp này còn góp phần giảm áp lực và trách nhiệm cho các cán bộ trung gian trong quá trình thu nộp.

Trên cơ sở các hướng giải quyết nêu trên, đề tài hướng đến việc xây dựng một cổng thông tin điện tử tập trung dành cho hoạt động Đoàn -- Hội, đóng vai trò là nền tảng quản lý thống nhất cho dữ liệu người dùng, hoạt động và tài chính. Đặc biệt, việc tích hợp công nghệ trí tuệ nhân tạo thông qua AI Chatbot hỗ trợ hỏi đáp tự động được kỳ vọng sẽ giúp sinh viên tiếp cận thông tin nhanh chóng, giảm tải công việc cho cán bộ và nâng cao hiệu quả vận hành của hệ thống.

